\section{The Deception of the New Age}

\begin{quotex}
Because of time constraints, \textbf{Alexander Shepard} was unable to post this piece himself, but has given us permission to reproduce his notes on the New Age as counter tradition. \textbf{Rene Guenon} writes of it:

\begin{quotex}

As for the counter tradition, we can still only see the preliminary signs of it, in the form of all things that are striving to become counterfeits in one way or another of the traditional idea itself.
\end{quotex}


The movements Mr. Shepard mentions all use the language of Tradition and even claim to be based on Traditional ideas and texts. But ``tradition'' means ``handed down'' and the revelation that was ``handed down'' is an ``inexhaustible deposit of wisdom which renews itself each time it is actualized by a new experience'' (Guido de Giorgio). But this is precisely what the New Age movements reject. That deposit never suffices for them, so they continually seek out novelty, either in the form of revolutionary social and political movements or in scientistic ideas such as Darwinism, Quantum Physics, and various psychological theories.

Men coming of age now have a unique opportunity. The texts of the founding fathers of Tradition are readily available, so there is no need to waste time with New Age ideas. That was not always so; in the next few days, I'll pass on my own experience with the New Age. We can only hope that this opportunity in time is not squandered. The fear is that what comes too easily is too little appreciated; there is something to be said about the need to for personal struggle. We conclude with a warning from Guenon, followed by Me. Shepard's text.

\begin{quotex}
Something much more formidable is in preparation for a future considered by some to be near, and the growing rapidity of the succession of events today is an indication of its proximity.
\end{quotex}

\end{quotex}


I thought I should offer my critique and perspective as to why ``New age'' thought, generally defined as any number of western ``paths'' arising in the west in the aftermath of the so-called ``Enlightenment'' is not a valid path. These spiritual movements include, but are certainly not limited to, Theosophy, the Bahia faith, Scientology, Wicca, Raelianism, Meher Baba's cult, The ``Universal Sufi'' cult, The Followers of Osho, the list goes on. Also included are those who are so found of calling themselves ``Spiritual but not religious''. There is particularly disturbing trend I have noticed among some Traditionalists of this generation to try to synthesize the Traditionalism and New Age. Books which come to mind includes a book called ``The Return of the Perennial Philosophy'' which is an account of Western Esoterism, and specifically criticizes Guenon for ``neglecting'' theosophy. Another philosopher who advocates such views is ``Joscelyn Godwin'' who in his book ``The Golden Thread'' seems to indicate a respect for the Theosophists and the new age movement in general. While he does not go out of his way to make a claim about the validity of the Theosophists, his works provide the basis for such a philosophy. While I am not saying that such people should be ignored, one should at least be cautious of them and be mindful of the advice of the 1st generation Traditionalists in deciphering True Mysticism from Sophistry.

\begin{enumerate}


\item ``New age thought'' claims to rest on the fulfillment of the Self. Yet they confuse the Psychic and the Spiritual. The Spiritual, or Soul, transcends the mind, and there can be no matter of self-realization. The Soul exists as a reflection of God. Any path not founded on surrendering their will and worshipping the Absolute will find no spiritual awakening. Treating the ``Ego'' or the ``individual self'' as the object of care is like trying to support and cancerous Tumor in the brain blocking rational thought. The correct approach is to DESTROY the tumor. This is the object of REAL mysticism, to slay this stain called the ``self''. As Al-Hallaj (ra) said ``I look for myself and the self was gone''

\item Many ``New agers'' seek an alternative to Abrahamic Spirituality. As a result, they openly scorn monotheism. Instead, they embrace ``animism'' or ``polytheism''. These two conceptions never existed; the notion of an ``Evolutionary'' thesis of religion was a fabrication of the western colonialists in the ``Age of Imperialism''. In reality, all Religion is Monotheistic, which is why it was easy for Africans and Native Americans to embrace Islam and Christianity, for the basic tenets of a Single all-powerful God, an afterlife, a fall from Grace, were already essential tenants of their religious beliefs. The New Agers who call themselves ``animists'' or ``polytheists'' are creating a gross-perversion and distortion of indigenous beliefs. Julius Evola made this case against the Neo-pagans claiming to practice European ``Paganism'', specially attacking their libertinism and ``nature worship''

\item On this note, The New Agers are often ignorant of Christian Mysticism. In their question for direct spirituality they oddly ignore clear examples of Self-realization within the Christian Tradition. One of the most clear examples being the great Meister Eckhart, whose entire Theology rested on demonstrating on how the Universe was an illusion and how the only thing Real was God, and that Christ was the ``Son of God'' in-so much as he was the most perfect Human, completing the levels of spiritual ascension, and providing an archetype for all human kind to reject their created natures and become One with God.

\item On this note, so-called ``New Agers'' are, like religious fundamentalists ignorant of the distinction between the esoteric and the exoteric. There is no reaching the Supreme Realization without rigorous bodily discipline. Even in the case of Primitive society, the Shaman had to subject himself to years and years of fasting, prayer, and seclusion. A great many trials and tribulations which the ``new ager'' could never survive. In fact, many of the most vocal critics of new age have been natives. In light of the Islamic tradition, the noble discipline of Tasawuff has been bastardized, with ``Neo-Sufis'' shamelessly claiming the mantel of ``Sufi'' without practicing the Shariah. They desecrate the graves of the sages of Islam by claiming absurdities like ``Sufism is different than Islam''. Authentic Spirituality has ALWAYS demanded the rigorous practice of the exoteric, and history, with rare exceptions, has demonstrated this. In order to transcend to elevate the spirit, we must conquer the body.
\end{enumerate}



\flrightit{Posted on 2012-06-19 by Admin}

\begin{center}* * *\end{center}

\begin{footnotesize}
\begin{sffamily}
\texttt{Avery on 2012-06-19 at 22:04 said:}

Guenon didn't neglect Theosophy. He wrote an entire book about it, which I think Mr. Shepard could learn a great deal from if he wants to learn more about these pseudo-traditional movements. However, I personally found this book wanting. What it should be doing is drawing a distinction between New Age methods of finding ``spiritual truth'' and Traditional methods. Instead he frequently fails to justify the theoretical basis for his attacks, it really just reads like a long, unedited Internet rant by a disillusioned ex-member (he was not a Theosophist but worked in similar occult groups). Evola never wrote anything so whiny about Fascism.

\hfill

\texttt{golgonooza on 2012-06-20 at 06:25 said:}

Surely you would want to support the bold claim that ``all Religion is Monotheistic'' with something a bit more substantial than the fact that Africans were able to accept Christianity easily. I think the reality is far more complex -- maybe the words monotheistic and polytheistic don't really do justice to it. Are Hindus monotheistic? Are they polytheistic? Aside from which school of thought you look at there would be problems with saying they were either.

\hfill

\texttt{Cologero on 2012-06-20 at 07:14 said:}

With all due respect, Avery, it is a common tactic of the modern mind to dismiss ideas based on the alleged psychological of the writer. You don't know if Guenon was disillusioned or not; perhaps he was trying to protects others from the initial illusion. He did, in fact, object to specific doctrines in Theosophy. Evola reprised the same arguments in his ``mask and face of contemporary spirituality''. If they misrepresented or misunderstood those doctrines, then specific examples would be of interest.

Evola wrote an entire book, critical of Fascsim'', ``Fascism seen from the right''. Do you not consider that ``whiny''? In any event, the only question is whether it is true or false.

\hfill

\texttt{Cologero on 2012-06-20 at 07:25 said:}

I agree, golgonooza, that alleged empirical facts about Africans cannot be used to \emph{prove} a principle, although facts may be helpful in illustrating it. Nevertheless, in the metaphysical sense, orthodox Traditions are monotheistic, that is, they agree there is an ultimate unifying principle. A literal polytheism implies there are multiple, independent, irreconcilable, ultimate principles.

\textit{THE DOCTRINE OF DIVINE UNITY IN HELLENIC TRADITION}\footnote{\url{http://www.eurasia-rivista.org/the-doctrine-of-divine-unity-in-hellenic-tradition/12164/}} may be of interest in this regard.

\hfill

\texttt{Avery Morrow on 2012-06-21 at 10:34 said:}

Cologero, I was not sure what to think of your reply on first read, but I now realize you are absolutely right and I will have to apologize for my bad judgment. Psychoanalysis of an author's mood and motivation is an anti-Traditional method of criticism and is actually reflective of the sort of cliquish bullying that is falsely passed off as critique on the Internet these days.

\hfill

\end{sffamily}
\end{footnotesize}

